
\subsubsection{Spurious Correlation Debiasing}

Neural networks exploit spurious correlations---features correlated with labels in training data but not causally informative---resulting in poor worst-group accuracy when spurious features misalign with labels at test time. Waterbirds exemplifies this: models trained with 95\% background-label correlation achieve 97\% accuracy on majority groups but collapse to 40\% on minority groups. CelebA exhibits gender-attribute spurious correlations, and CMNIST provides a controlled testbed with color-label correlation.

\textbf{Group-based methods} improve robustness by upweighting minority groups during training. GroupDRO performs distributionally robust optimization over worst-case group risk, while Invariant Risk Minimization learns predictors invariant across environments. However, these methods require group annotations unavailable in many deployment scenarios.

\textbf{Reweighting methods} address spurious correlations without group labels by exploiting learning dynamics. Just Train Twice (JTT) trains a biased ERM model, identifies hard examples (those the biased model misclassifies), then retrains with upweighted hard examples. Learning from Failure extends this by automatically discovering groups via prediction disagreement across training checkpoints. Both methods hypothesize that spurious features are learned earlier than core features, enabling hard example identification to isolate core-feature-dependent examples. However, neither paper measures temporal ordering at the gradient level---this work provides the first direct validation via per-epoch gradient norm tracking on ablated feature types.

\subsubsection{Learning Dynamics and Example Difficulty}

Prior work analyzes learning dynamics at the example level, tracking which samples are forgotten or misclassified during training. Toneva et al. introduce the forgetting metric---the number of times an example's prediction flips from correct to incorrect across epochs---showing that ''unforgettable'' examples are learned early and stably, while ''forgettable'' examples exhibit oscillating predictions. Swayamdipta et al. propose data maps that cluster examples by confidence and variability, identifying easy vs hard examples.

These example-level analyses inform data pruning and curriculum learning but do not isolate feature-level learning dynamics. An example may be hard because it contains only core features, or because core and spurious features conflict---example-level hardness does not distinguish feature types. This work adapts forgetting analysis to the feature level: spurious-only and core-only network variants are trained via ablation, then forgetting rates per feature type are measured. Spurious features show $F_{\text{spurious}} = 2.35$ vs $F_{\text{core}} = 4.82$ events/sample, confirming that spurious features converge to more stable predictions---a finding consistent with simpler decision boundaries but not observable from example-level metrics alone.

\textbf{Loss landscape and sharpness.} Sharpness-Aware Minimization (SAM) improves generalization by minimizing both loss value and loss sharpness (largest Hessian eigenvalue). Keskar et al. show that large-batch training converges to sharp minima with poor generalization, while small-batch training finds flat minima. SAM's connection to spurious correlation robustness remains underexplored. Layer-wise neuron correlation analysis provides preliminary evidence: early layers show higher neuron-spurious correlation than late layers, consistent with spurious features being localized in shallow, potentially sharper regions of the loss landscape. However, Hessian eigenspectrum is not directly measured, leaving sharpness-temporal gap connections for future work.

\subsubsection{Optimization Implicit Bias}

Gradient descent exhibits implicit bias toward specific solutions even without explicit regularization. Soudry et al. prove that gradient descent on linearly separable data converges to the max-margin solution in the direction of weights, even with zero regularization. Gunasekar et al. extend this to matrix factorization, showing implicit bias toward low nuclear norm. Neyshabur et al. analyze implicit bias in deep networks via PAC-Bayes bounds, suggesting that SGD favors solutions with small effective capacity.

These theoretical results establish that gradient descent prefers simpler solutions, but ''simplicity'' is defined in parameter space (margin, norm, rank) rather than feature space. This work's temporal gap measurement bridges this gap: if spurious features provide simpler decision boundaries than core features, implicit bias should manifest as earlier convergence for spurious features. Ablation training isolates feature types, enabling direct measurement: spurious-only training converges at $E_s=13$, while core-only training converges at $E_c=17$, yielding $\Delta=4$ epochs temporal gap. This empirical validation on realistic spurious benchmarks extends implicit bias theory from simplified linear settings to hierarchical deep networks on vision tasks.

\textbf{Architectural inductive bias.} CNNs process images via local receptive fields that expand hierarchically (early layers capture low-level features like edges and color, late layers capture high-level semantics like object parts). Vision Transformers (ViTs) use global self-attention from the first layer, potentially processing low- and high-level features more uniformly. CNNs may exhibit larger temporal gaps due to hierarchical layer-wise feature processing amplifying the simplicity bias toward early-layer spurious features. Architectural comparison experiments tested this on CMNIST, though initial results were inconclusive due to convergence threshold design issues (both architectures converged within 1-2 epochs at 70\% accuracy threshold, providing insufficient temporal window). Future work with higher thresholds (85-90\%) and longer training (100 epochs) is needed to validate architectural modulation.
